\section{The cross-provider protection level and the integrity envelope}
\label{sec:cmpl}

We now turn the blind magnitude of Section~\ref{sec:anchor} into an integrity bound.
Every quantity in this section is \Modelled{} and tagged \modtag{}: the constellation
geometry, the common-mode covariance, the cross-provider protection level (CMPL), the
representative ARAIM horizontal protection level (HPL), and the total envelope are all
representative---the ``providers'' are ephemerides mapped through a modelled
constellation, and \emph{no two lunar-PNT providers fly}. The section's purpose is to
state the load-bearing claim in its correct form and to exhibit the envelope's
structure, not to certify a number.

\subsection{Why a new term is needed: the load-bearing claim, stated correctly}
\label{sec:loadbearing}

The tempting claim---``ARAIM is provably blind to common-mode''---is false, and a
reviewer is right to reject it. The distinction is between the residual test and the
protection level.

\begin{itemize}
  \item The RAIM/ARAIM \emph{residual test} is provably blind to $\rng(\Gmat)$
  (Proposition~\ref{prop:blind}, Lemma~\ref{lem:rigid}). A common-mode frame/ephemeris
  error triggers no fault flag.

  \item The ARAIM \emph{protection level} does \emph{not} rely on the residual test to
  catch this error. The multiple-hypothesis solution-separation (MHSS) protection level
  bounds the position error under the \emph{fault-free} (nominal, $H_0$) hypothesis
  through the user range accuracy overbound $\sigma_{\mathrm{URA}}$ and the nominal bias
  $b_{\mathrm{nom}}$ declared in the integrity support message (ISM)
  \citep{blanch2015,wgc2016araim,decleene2000}. A residual-invisible error is therefore
  still budgeted---up to whatever the ISM declares. Within a single provider, an
  ephemeris/frame bias that the residual test cannot see is covered by
  $\sigma_{\mathrm{URA}}$ and $b_{\mathrm{nom}}$, provided the ISM overbound is honest
  \citep{montenbruck2018sisre}.
\end{itemize}

The genuinely uncovered term is specific to the \emph{multi-provider} case. When a user
fuses, or hands off between, two \emph{independent} lunar providers, their mutual
frame/ephemeris common-mode is not contained in \emph{either} provider's per-provider
URA. Each provider's $\sigma_{\mathrm{URA}}$ overbounds \emph{its own} signal-in-space
error against \emph{its own} realized frame and ephemeris; it says nothing about the
offset between two independent realizations. The measured inter-ephemeris
common-mode---$\approx\!2\,$m between DE440, INPOP21a, and EPM2021
(Section~\ref{sec:anchor}, and the interoperability floor of \citep{baweja2026interop})%
---is exactly such a cross-provider term. It is residual-invisible (Lemma~\ref{lem:rigid})
\emph{and} outside any single provider's ISM budget. That is the gap the CMPL fills.

\subsection{The common-mode covariance and the CMPL}
\label{sec:cmpldef}

We model the cross-provider common shift as a zero-mean Gaussian with an isotropic
$3$-D covariance whose total one-sigma magnitude is fixed to the Part-B floor,
\begin{equation}
  \Sig_{\mathrm{cm}} \;=\; \frac{\sigma_{\mathrm{cm}}^2}{3}\,\bm I_3,
  \qquad
  \sigma_{\mathrm{cm}} \;=\; \|\ds\| \;=\; \SI{2.0000}{m}\ \text{(1-}\sigma\text{)},
  \label{eq:cmcov}
\end{equation}
sized to the $\approx\!2\,$m inter-ephemeris floor measured in Section~\ref{sec:anchor}
and in the companion interoperability paper \citep{baweja2026interop}, which is the
input that feeds $\Sig_{\mathrm{cm}}$. Because a common shift maps (Lemma~\ref{lem:rigid})
directly into user position error, the blind position covariance is $\Sig_{\mathrm{cm}}$
itself. Projecting onto the horizontal plane gives a horizontal one-sigma
\begin{equation}
  \sigma_H \;=\; \sqrt{\tfrac{2}{3}}\,\sigma_{\mathrm{cm}} \;=\; \SI{1.6330}{m},
  \label{eq:sigmaH}
\end{equation}
(two of three isotropic axes are horizontal), and the cross-provider protection level is
the standard $2$-D horizontal Rayleigh inflation of this one-sigma,
\begin{equation}
  \mathrm{CMPL} \;=\; K_H\,\sigma_H \;=\; 6.0\times\SI{1.6330}{m} \;=\; \SI{9.7980}{m},
  \qquad K_H = 6.0,
  \label{eq:cmpl}
\end{equation}
with $K_H=6.0$ the representative horizontal integrity multiplier (the $2$-D Rayleigh
factor used at civil-aviation integrity risks in ARAIM \citep{blanch2015,wgc2016araim}).
We label \eqref{eq:cmpl} plainly: the CMPL value is an \emph{input-echo mechanism
illustration}, not a finding. It restates the $\approx\!2\,$m modelled input through a
fixed multiplier; its role is to show \emph{how} a cross-provider term enters an
integrity budget and with what sign and magnitude, not to certify $\SI{9.80}{m}$ as a
lunar CMPL.

\subsection{The integrity envelope, and why the sum is a triangle bound}
\label{sec:envelope}

We combine the CMPL with a representative single-provider ARAIM horizontal protection
level. Taking $\mathrm{HPL}_{\mathrm{ARAIM}}=\SI{15.0}{m}$ (representative and
illustrative---\emph{not} \Validated{}; a placeholder for whatever a real lunar ARAIM
configuration would compute), the total cross-provider integrity envelope is
\begin{equation}
  \mathrm{HPL}_{\mathrm{total}}
  \;=\; \mathrm{HPL}_{\mathrm{ARAIM}} + \mathrm{CMPL}
  \;=\; \SI{15.0}{m} + \SI{9.7980}{m}
  \;=\; \SI{24.7980}{m}.
  \label{eq:hpltotal}
\end{equation}
Table~\ref{tab:envelope} collects the terms. Two properties of \eqref{eq:hpltotal} are
essential and both are stated as \Modelled{} guidance.

\begin{table}[t]
\centering
\caption{The cross-provider integrity envelope (\Modelled{}, \modtag{}). The
common-mode one-sigma $\sigma_{\mathrm{cm}}=\SI{2}{m}$ is sized to the real
inter-ephemeris floor (Section~\ref{sec:anchor}); the CMPL is an input-echo mechanism
illustration, not a finding. The total is a \emph{conservative} triangle-inequality sum:
a root-sum-square ($\sqrt{15^2+9.798^2}=\SI{17.92}{m}$) would under-bound and overclaim
integrity, because the ARAIM and common-mode contributions can align in the position
domain. The envelope is valid only where the common-mode is not already absorbed into
$\sigma_{\mathrm{URA}}$ (Section~\ref{sec:conditionality}).}
\label{tab:envelope}
\small
\begin{tabular}{@{}lcl@{}}
\toprule
\textbf{Term} & \textbf{Value (m)} & \textbf{Status} \\
\midrule
$\sigma_{\mathrm{cm}}=\|\ds\|$ (common-mode 1-$\sigma$) & $2.0000$ & \Modelled{} input, sized to real floor \\
$\sigma_H=\sqrt{2/3}\,\sigma_{\mathrm{cm}}$ (horizontal 1-$\sigma$)   & $1.6330$ & \Modelled{} \\
$\mathrm{CMPL}=K_H\sigma_H$, $K_H=6.0$    & $9.7980$ & \Modelled{} (input-echo illustration) \\
$\mathrm{HPL}_{\mathrm{ARAIM}}$ (representative)          & $15.0000$ & \Modelled{}, \emph{not} \Validated{} \\
$\mathrm{HPL}_{\mathrm{total}}=\mathrm{HPL}_{\mathrm{ARAIM}}+\mathrm{CMPL}$ & $24.7980$ & \Modelled{} (conservative triangle sum) \\
\midrule
(root-sum-square, for contrast---\emph{under-bounds}) & $17.92$ & rejected: overclaims integrity \\
\bottomrule
\end{tabular}
\end{table}

\paragraph{The sum is a triangle bound, not a root-sum-square.}
Integrity budgets normally combine independent error sources in quadrature. Here that
would be wrong. The single-provider ARAIM error and the cross-provider common-mode are
\emph{not} guaranteed independent in the position domain: their horizontal
contributions can align, so the worst-case combined displacement is the arithmetic sum,
not the root-sum-square. Using $\sqrt{\mathrm{HPL}_{\mathrm{ARAIM}}^2+\mathrm{CMPL}^2}
=\SI{17.92}{m}$ would under-bound the true worst case of $\SI{24.80}{m}$ and
\emph{overclaim} integrity---the cardinal sin of an integrity monitor. The
triangle-inequality sum \eqref{eq:hpltotal} is the conservative choice and the correct
one for a bound.

\subsection{The anti-double-counting conditionality}
\label{sec:conditionality}

The envelope \eqref{eq:hpltotal} is valid \emph{only where the common-mode is not
already absorbed into $\sigma_{\mathrm{URA}}$}. This condition is not a footnote; it is
the boundary of the CMPL's validity.

\begin{itemize}
  \item If each provider's ISM declares an \emph{un-inflated} $\sigma_{\mathrm{URA}}$
  that covers only its own signal-in-space error against its own realization---the
  normal case---then the inter-provider common-mode is genuinely uncovered, and adding
  the CMPL is necessary. The envelope \eqref{eq:hpltotal} then restores a valid bound.

  \item If, instead, a provider inflates $\sigma_{\mathrm{URA}}$ (or $b_{\mathrm{nom}}$)
  to \emph{also} cover the cross-provider disagreement, then the common-mode is already
  inside the ARAIM fault-free term, and adding the CMPL \emph{double-counts}: the true
  bound is already $\mathrm{HPL}_{\mathrm{ARAIM}}$ alone, and
  $\mathrm{HPL}_{\mathrm{ARAIM}}+\mathrm{CMPL}$ is merely conservative-to-the-point-of-%
  waste, not wrong, but no longer tight.
\end{itemize}

Hence the precise statement:
$\mathrm{HPL}_{\mathrm{total}}=\mathrm{HPL}_{\mathrm{ARAIM}}+\mathrm{CMPL}$ is the valid
bound for the \emph{un-inflated-URA} case---the case where a provider budgets only its
own realization and the cross-provider term is left to a multi-provider integrity layer.
The design consequence for a multi-provider architecture is a clean choice: either
mandate that each provider's ISM overbound explicitly \emph{includes} the permissible
inter-provider common-mode (folding CMPL into $\sigma_{\mathrm{URA}}$ upstream), or
carry a separate CMPL term downstream at the fusion point. What must not happen is
\emph{neither}---leaving the cross-provider common-mode in no budget at all---which is
exactly the gap Lemma~\ref{lem:rigid} and Section~\ref{sec:anchor} expose: it is
residual-invisible \emph{and}, absent an explicit provision, outside every
per-provider URA.
